% mnras_template.tex 
%
% LaTeX template for creating an MNRAS paper
%
% v3.3 released April 2024
% (version numbers match those of mnras.cls)
%
% Copyright (C) Royal Astronomical Society 2015
% Authors:
% Keith T. Smith (Royal Astronomical Society)

% Change log
%
% v3.3 April 2024
%   Updated \pubyear to print the current year automatically
% v3.2 July 2023
%	Updated guidance on use of amssymb package
% v3.0 May 2015
%    Renamed to match the new package name
%    Version number matches mnras.cls
%    A few minor tweaks to wording
% v1.0 September 2013
%    Beta testing only - never publicly released
%    First version: a simple (ish) template for creating an MNRAS paper

%%%%%%%%%%%%%%%%%%%%%%%%%%%%%%%%%%%%%%%%%%%%%%%%%%
% Basic setup. Most papers should leave these options alone.
\RequirePackage{fix-cm}
\documentclass[fleqn,usenatbib]{mnras}
% MNRAS is set in Times font. If you don't have this installed (most LaTeX
% installations will be fine) or prefer the old Computer Modern fonts, comment
% out the following line
\usepackage{newtxtext,newtxmath}
% Depending on your LaTeX fonts installation, you might get better results with one of these:
%\usepackage{mathptmx}
%\usepackage{txfonts}

% Use vector fonts, so it zooms properly in on-screen viewing software
% Don't change these lines unless you know what you are doing
\usepackage[T1]{fontenc}

% Allow "Thomas van Noord" and "Simon de Laguarde" and alike to be sorted by "N" and "L" etc. in the bibliography.
% Write the name in the bibliography as "\VAN{Noord}{Van}{van} Noord, Thomas"
\DeclareRobustCommand{\VAN}[3]{#2}
\let\VANthebibliography\thebibliography
\def\thebibliography{\DeclareRobustCommand{\VAN}[3]{##3}\VANthebibliography}

% Only include extra packages if you really need them. Avoid using amssymb if newtxmath is enabled, as these packages can cause conflicts. newtxmatch covers the same math symbols while producing a consistent Times New Roman font. Common packages are:
\usepackage{graphicx}	% Including figure files
\usepackage{amsmath}	% Advanced maths commands

%%%%%%%%%%%%%%%%%%% TITLE PAGE %%%%%%%%%%%%%%%%%%%

% Title of the paper, and the short title which is used in the headers.
% Keep the title short and informative.
\title[Short title, max. 45 characters]{A Phenomenological Acceleration Law from Galaxy Dynamics Consistent with Type Ia Supernova Distances}

% The list of authors, and the short list which is used in the headers.
% If you need two or more lines of authors, add an extra line using \newauthor
\author[Donald Airey]{
Donald Airey,$^{1}$\thanks{E-mail: don@airey.us}
\\
% List of institutions
$^{1}$Worcester Polytechnic Institute, 100 Institution Road, Worcester, MA 01609, USA\\
}

% These dates will be filled out by the publisher
\date{Accepted XXX. Received YYY; in original form ZZZ}

% Prints the current year, for the copyright statements etc. To achieve a fixed year, replace the expression with a number. 
\pubyear{\the\year{2026}}

% Don't change these lines
\begin{document}
\label{firstpage}
\pagerange{\pageref{firstpage}--\pageref{lastpage}}
\maketitle

% Abstract of the paper
\begin{abstract}
We test whether a single empirically determined acceleration scale can link galaxy dynamics and cosmological distance measurements. A phenomenological acceleration law is fit directly to galaxy kinematic data from the SPARC sample, yielding a fixed acceleration scale and reproducing the baryonic Tully--Fisher relation as a derived consequence. This scale defines an analytic, closed-form distance--redshift relation, which is evaluated against the Pantheon+SH0ES Type~Ia supernova dataset in a predictive setting.

The resulting relation predicts Type~Ia supernova distances with a significantly lower covariance-weighted $\chi^2$ than a spatially flat FLRW model with Planck 2018 parameters under identical, out-of-sample evaluation conditions, providing a compact and computationally efficient representation of the observed distance--redshift relation in this empirical test.
\end{abstract}

% Select between one and six entries from the list of approved keywords.
% Don't make up new ones.
\begin{keywords}
cosmology: observations -- distance scale -- supernovae: general -- galaxies: kinematics and dynamics
\end{keywords}

%%%%%%%%%%%%%%%%% BODY OF PAPER %%%%%%%%%%%%%%%%%%

\section{Introduction}

Type~Ia supernovae (SNe~Ia) provide one of the most direct observational probes of the cosmic expansion history through the luminosity distance--redshift relation. Large compilations such as Pantheon+~\cite{Brout2022}, building on earlier compilations~\cite{Scolnic2018}, and constructed using standard light-curve fitting techniques such as SALT2~\cite{Guy2007}, enable high-precision tests of cosmological models using covariance-weighted statistical methods. In the standard framework, Type~Ia supernova observations are interpreted within a parameterized expansion history specified by the Hubble constant and the relative contributions of matter, radiation, and dark energy, with spatial flatness often assumed. These parameters determine the redshift dependence of the expansion rate $H(z)$ and, through it, the luminosity distance.

An alternative approach is to consider constrained analytic forms for the distance--redshift relation and to evaluate their performance directly against the data. In this work, we examine a model in which the distance relation is parameterized by a small set of kinematic quantities, one of which is an acceleration scale determined independently from galaxy dynamics. This approach enables a direct test of whether a single empirically determined scale can consistently describe both galactic and cosmological observables.

The acceleration scale is obtained by fitting the phenomenological acceleration law directly to galaxy kinematic data from the SPARC sample, yielding a fixed value that is not adjusted in the supernova analysis. The remaining parameter is calibrated using the Cepheid subset of the Pantheon+ sample, and the resulting model is evaluated on the out-of-sample supernovae using the full covariance matrix.

The goal of this paper is to assess whether this constrained analytic relation provides a quantitatively consistent description of the observed SNe~Ia distance moduli across redshift under identical calibration conditions. No assumption is made regarding the underlying gravitational or cosmological theory, and the analysis is restricted to the observables considered here.

\section{Phenomenological Acceleration Law}

The purpose of this work is to test whether a single empirically determined acceleration scale can link two otherwise distinct observational regimes: circular motion in rotationally supported galaxies and the luminosity distances of Type~Ia supernovae. We define a phenomenological acceleration law for circular orbits, determine the acceleration scale $A$ from the baryonic Tully--Fisher relation, and then hold it fixed in the supernova analysis. No underlying gravitational theory is assumed. This separation is essential: the acceleration scale entering the distance--redshift relation is not fitted to the supernova sample.

\subsection{Circular Orbits}

We consider circular motion in an acceleration field consisting of a Newtonian component and a constant inward term $A$. For a test body in a circular orbit of radius $r$ about an enclosed baryonic mass $M_b(r)$, the radial acceleration is
\begin{equation}
a(r) = -A - \frac{G M_b(r)}{r^2}.
\end{equation}

The condition for circular motion requires that this radial acceleration provide the centripetal acceleration,
\begin{equation}
-\frac{v^2}{r} = -A - \frac{G M_b(r)}{r^2}.
\end{equation}
Multiplying by $r$ gives
\begin{equation}
v^2 = A r + \frac{G M_b(r)}{r},
\end{equation}
and solving for the enclosed baryonic mass yields
\begin{equation}
M_b(r) = \frac{r\left(v^2 - A r\right)}{G}.
\label{eq:massradius}
\end{equation}

\subsection{Circular-Orbit Solutions}

Equation~(\ref{eq:massradius}) defines a two-parameter family of circular-orbit solutions relating baryonic mass $M$, orbital radius $r$, and orbital speed $v$.

For fixed orbital speed $v$, the mass function $M(r)$ is a concave quadratic in $r$. Differentiating gives
\begin{equation}
\frac{dM}{dr} = \frac{v^2 - 2 A r}{G} .
\end{equation}
The mass therefore reaches a maximum at
\begin{equation}
r_{\max} = \frac{v^2}{2A} .
\end{equation}
Substituting this radius into equation~(\ref{eq:massradius}) yields the maximum baryonic mass compatible with a circular orbit at speed $v$,
\begin{equation}
M_{\max}(v) = \frac{v^4}{4 A G}.
\label{eq:btfr_envelope}
\end{equation}

The surface $M(v,r)$ defines a family of circular-orbit solutions, as shown in Figure~\ref{fig:fundamentalplane}, with the ridge $M_{\max}(v)$ giving the upper envelope of baryonic mass for a given orbital speed. Galaxies dominated by circular rotation are expected to saturate this boundary and follow a characteristic mass--velocity relation.

\begin{figure}
\centering
\includegraphics[width=0.8\linewidth]{FundamentalPlane.pdf}
\caption{Circular-orbit solution surface defined by equation~(\ref{eq:massradius}). The baryonic mass surface $M(v,r)$ is shown as a function of orbital speed $v$ and orbital radius $r$ for an illustrative acceleration scale $A = 10^{-11}\,\mathrm{m\,s^{-2}}$. The red ridge marks the locus $M_{\max}(v)$ corresponding to the maximum baryonic mass permitted at each orbital speed. Galaxies populating this ridge reproduce the baryonic Tully--Fisher relation.}
\label{fig:fundamentalplane}
\end{figure}

\subsection{Baryonic Tully--Fisher Relation}

$M_{\max}(v)$ defines the scaling $M \propto v^4$, consistent with the observed baryonic Tully--Fisher relation in~\cite{McGaugh2000}. The normalization of this relation is set by the acceleration scale $A$. A characteristic acceleration scale in galaxy dynamics has previously been emphasized in frameworks such as MOND~\cite{Milgrom1983}.

This parameter is determined empirically by fitting the circular-orbit envelope (equation~(\ref{eq:btfr_envelope})) to the SPARC galaxy sample of McGaugh, Lelli and Schombert~\cite{McGaugh2016}, using the assumptions $\Upsilon_\star=0.5$, $f_{\mathrm{gas}}=1.33$, zero intrinsic mass-to-light uncertainty, and the quality cut $\mathrm{Qual}\leq 3$. These rotationally supported disk systems are well suited to this analysis, as their kinematics closely follow circular-orbit behavior.

For each galaxy the rotation curve provides an asymptotic circular velocity $v$, while the baryonic mass $M_b$ is computed as $M_b=\Upsilon_\star L_{3.6}+f_{gas}\,M_{\mathrm{HI}}$. The value of $A$ is determined by minimizing the $\chi^2$ statistic between the theoretical envelope and the observed $(M_b,v)$ pairs.

The best-fit acceleration scale is
\begin{equation}
A = 3.7\times10^{-11}\,\mathrm{m\,s^{-2}} .
\end{equation}
The observed galaxies follow the predicted $M\propto v^4$ scaling, as shown in Figure~\ref{fig:btfrchart}. In this construction, the baryonic Tully--Fisher relation is not imposed but follows directly from the acceleration law, with the constant acceleration scale setting its normalization.

\begin{figure}
\centering
\includegraphics[width=0.75\linewidth]{BTFRChart.pdf}
\caption{Log--log plot of baryonic mass $M_b$ versus outer rotation speed $v$ for the SPARC galaxy sample. Black circles show the observed galaxies. The solid red line shows the circular-orbit envelope evaluated at the best-fit acceleration, demonstrating consistency with the observed $M \propto v^4$ scaling.}
\label{fig:btfrchart}
\end{figure}

\subsection{Distance--Redshift Relation}

With the acceleration scale fixed empirically from galaxy dynamics, we now construct a phenomenological analytic form for the comoving distance,
\begin{equation}
D_C(z) = \frac{t_o\left(2V_0 - A t_o\right) z}{2 + z},
\end{equation}
where $t_o$ and $V_0$ set the temporal and velocity scales of the relation, respectively. The luminosity distance follows as
\begin{equation}
D_L(z) = (1+z)D_C(z).
\end{equation}
This closed-form expression differs from the standard FLRW integral representation and provides an alternative distance--redshift relation for direct comparison with observations.

\subsection{Distance Modulus}

Type~Ia supernova observations are reported in terms of the distance modulus~\cite{Tripp1998}, which relates the observed luminosity distance to the measured brightness of each supernova,
\begin{equation}
\mu(z) = 5\log_{10}\!\left(\frac{D_L(z)}{10\,\mathrm{pc}}\right).
\end{equation}

\subsection{Kinematic Parameters}

We treat the parameters in the analytic distance--redshift relation as kinematic quantities that set the temporal, velocity, and acceleration scales of the model, without assuming an underlying dynamical theory.

A boundary condition is imposed at the observation time $t_o$ by requiring that the parameter combination reproduce the locally measured speed of light,
\begin{equation}
c = V_0 - A t_o .
\label{eq:co_definition}
\end{equation}
This condition fixes the integration constant $V_0$ in the analytic distance relation. Rearranging,
\begin{equation}
V_0 = c + A t_o .
\label{eq:velocity_boundary_condition}
\end{equation}

With the acceleration scale $A$ fixed from the galaxy kinematic analysis, the Cepheid subset~\cite{Riess2022} is used to determine the temporal scale $t_o$, leaving a single free parameter. For each candidate value of $t_o$, theoretical distance moduli $\mu_{\mathrm{th}}(z)$ are evaluated and compared with the observed values $\mu_{\mathrm{obs}}(z)$, defining residuals
\begin{equation}
r_i = \mu_{\mathrm{obs}}(z_i) - \mu_{\mathrm{th}}(z_i).
\end{equation}
The best-fit value of $t_o$ is obtained by minimizing the covariance-weighted statistic
\begin{equation}
\chi^2 = r^{\mathrm T} C^{-1} r,
\end{equation}
where $C$ is the full Pantheon+ STAT+SYS covariance matrix. The resulting value of $t_o$ determines $V_0$, completing the kinematic parameter set. The fitted parameters are listed in Table~\ref{tab:model_parameters}.

\begin{table}
\centering
\begin{tabular}{|l|l|}
\hline
Parameter & Value \\
\hline
$A$ & $3.7\times10^{-11}\,\mathrm{m\,s^{-2}}$ \\
$V_0$ & $3.16\times10^{8}\,\mathrm{m\,s^{-1}}$ \\
$t_o$ & $4.51\times10^{17}\,\mathrm{s}$ \\
\hline
\end{tabular}
\caption{Kinematic parameters determined from galaxy kinematics and the Cepheid-calibrated Pantheon+ subset.}
\label{tab:model_parameters}
\end{table}

\subsection{Type Ia Supernovae}

\begin{figure}
\centering
\includegraphics[width=0.7\linewidth]{HubbleDiagram.pdf}
\caption{Pantheon+SH0ES distance moduli as a function of redshift $z$, compared with the analytic model prediction and the spatially flat FLRW model. Model parameters are fixed by BTFR and Cepheid calibrations, and the comparison is shown for the remaining supernovae. FLRW parameters are fixed to the Planck 2018 values. Residuals are shown for the analytic model (red circles) and FLRW (blue triangles).}
\label{fig:hubblechart}
\end{figure}

We compare the analytic distance--redshift relation with the Pantheon+SH0ES Type~Ia supernova dataset using covariance-weighted statistics. The resulting Hubble diagram and residuals are shown in Figure~\ref{fig:hubblechart}.

The Pantheon+SH0ES dataset~\cite{Brout2022} provides observed distance moduli $\mu_{\mathrm{obs}}(z)$ calibrated on an absolute scale via the SH0ES determination of the Type~Ia supernova absolute magnitude~\cite{Riess2022}. This calibration is applied uniformly to both models.

Each model predicts $D_L(z)$, which is mapped to $\mu_{\mathrm{th}}(z)$ using the definition of the distance modulus. We compare $\mu_{\mathrm{th}}(z)$ to $\mu_{\mathrm{obs}}(z)$ using the full Pantheon+ STAT+SYS covariance matrix.

We do not refit the FLRW model to the supernova sample, as this would convert the comparison into an in-sample descriptive fit rather than an out-of-sample predictive test. All models are therefore evaluated with parameters fixed independently of the supernovae used in the comparison. The analytic relation parameters are determined from galaxy kinematics and from Cepheid variable stars in the Pantheon+ dataset. This set of calibration records is then excluded from the evaluation, so the remaining supernovae provide an out-of-sample test of the predicted distance--redshift relation. The FLRW model is evaluated in a spatially flat form with parameters fixed to the Planck 2018 values, $\Omega_m = 0.315$, $\Omega_\Lambda = 1-\Omega_m$, and $H_0 = 67.4\,\mathrm{km\,s^{-1}\,Mpc^{-1}}$~\cite{Planck2018}.

The resulting covariance-weighted chi-square statistics are summarized in Table~\ref{tab:chi2_comparison}.

\begin{table}
\centering
\begin{tabular}{lccc}
\hline
Model & $\chi^2$ & $\nu$ & $\chi^2/\nu$ \\
\hline
Analytic model & 2726 & 1625 & 1.68 \\
FLRW (Planck 2018) & 4049 & 1625 & 2.49 \\
\hline
\end{tabular}
\caption{Covariance-weighted $\chi^2$ statistics for the analytic model and the spatially flat FLRW model evaluated on the Pantheon+SH0ES supernova sample.}
\label{tab:chi2_comparison}
\end{table}

The analytic relation yields a substantially lower covariance-weighted $\chi^2$ and reduced $\chi^2$ than the Planck 2018 FLRW model. The residuals as a function of redshift do not exhibit a clear systematic trend.

\section*{Data Availability}

The data underlying this article are publicly available. The SPARC galaxy sample can be accessed at \url{http://astroweb.cwru.edu/SPARC/}, and the Pantheon+SH0ES Type~Ia supernova dataset is available from the Pantheon+ data release at \url{https://github.com/PantheonPlusSH0ES/DataRelease} and the associated publications~\cite{Brout2022,Riess2022}. A computational notebook implementing the distance--redshift relation and observational analysis is available at \url{https://www.wolframcloud.com/obj/ed737c03-5baf-49f4-90e3-83b91b245c2a}. The notebook reproduces the results and figures presented in this work and is provided to enable independent verification of the calculations.

\section{Conclusion}

We test a phenomenological construction in which a single empirically determined acceleration scale links galaxy dynamics and cosmological distance measurements. The scale is fixed from galaxy kinematics, and a single additional parameter, calibrated on Cepheid variables, defines an analytic, closed-form distance--redshift relation.

When evaluated on the Pantheon+SH0ES supernova sample with parameters fixed in advance, the resulting relation yields a substantially lower covariance-weighted $\chi^2$ than a spatially flat $\Lambda$CDM model with Planck 2018 parameters, constituting an out-of-sample test of predictive performance.

The construction is directly falsifiable: with the acceleration scale and kinematic parameters fixed, it predicts a specific departure from Keplerian motion in bound systems such as the Solar System. This signal is absent from standard FLRW-based local dynamics without an additional local correction and is not a generic feature of MOND-like constructions. A robust non-detection would rule out the construction, while detection would provide a direct observational discriminator.

\bibliographystyle{mnras}
\bibliography{references}
\label{lastpage}

\end{document}